\chapter{What Is Inside an FPGA}
\label{ch:inside}

We have learned to design combinational circuits, flip-flops, registers, and
state machines and to check them in Logra. One question remains: what is such a
circuit made of inside the chip? An FPGA has no separate gates and no wires you
could re-solder. This chapter opens up the chip and shows how the blocks we drew
and described in Verilog fit into its resources.

If you took the lid off an FPGA and looked at the chip under a microscope,
you would see a neat ``chessboard'': identical logic blocks surrounded by
a network of wires, and input/output blocks along the edges. Let us go
through these parts one by one.

\begin{figure}[H]
\centering
\begin{tikzpicture}[scale=0.87]
  % Chip
  \fill[ChipBlack!8] (-0.6,-0.6) rectangle (12.6,8.6);
  \draw[very thick, FpgaDark] (-0.6,-0.6) rectangle (12.6,8.6);
  % Routing channels
  \foreach \x in {0,...,12} \draw[FpgaOrange!55, line width=0.7pt] (\x,0) -- (\x,8);
  \foreach \y in {0,...,8} \draw[FpgaOrange!55, line width=0.7pt] (0,\y) -- (12,\y);
  % Logic blocks / memory / DSP
  \foreach \x in {0,...,11} \foreach \y in {0,...,7} {
    \pgfmathtruncatemacro{\col}{\x}
    \ifnum\col=4
      \fill[FpgaPurple!65] (\x+0.15,\y+0.15) rectangle ++(0.7,0.7);
    \else\ifnum\col=8
      \fill[FpgaRed!65] (\x+0.15,\y+0.15) rectangle ++(0.7,0.7);
    \else
      \fill[FpgaBlue!75] (\x+0.15,\y+0.15) rectangle ++(0.7,0.7);
    \fi\fi
  }
  % Switches at the crossings
  \foreach \x in {0,...,12} \foreach \y in {0,...,8}
    \fill[FpgaOrange] (\x,\y) circle (0.06);
  % Input/output blocks
  \foreach \x in {0,...,11} {
    \fill[FpgaGreen] (\x+0.25,-0.5) rectangle ++(0.5,0.3);
    \fill[FpgaGreen] (\x+0.25,8.2) rectangle ++(0.5,0.3);
  }
  \foreach \y in {0,...,7} {
    \fill[FpgaGreen] (-0.5,\y+0.25) rectangle ++(0.3,0.5);
    \fill[FpgaGreen] (12.2,\y+0.25) rectangle ++(0.3,0.5);
  }
  % PLL
  \fill[yellow!80!orange] (-0.55,-0.55) rectangle (-0.05,-0.05);
  \fill[yellow!80!orange] (12.05,8.05) rectangle (12.55,8.55);
  % Legend
  \begin{scope}[xshift=13.3cm, yshift=7.4cm, every node/.style={font=\sffamily\small, anchor=west}]
    \fill[FpgaBlue!75] (0,0) rectangle (0.4,0.4); \node at (0.5,0.2) {CFU (logic)};
    \fill[FpgaPurple!65] (0,-0.8) rectangle (0.4,-0.4); \node at (0.5,-0.6) {BSRAM (memory)};
    \fill[FpgaRed!65] (0,-1.6) rectangle (0.4,-1.2); \node at (0.5,-1.4) {DSP (multipliers)};
    \fill[FpgaGreen] (0,-2.4) rectangle (0.4,-2.0); \node at (0.5,-2.2) {IOB (input/output)};
    \fill[yellow!80!orange] (0,-3.2) rectangle (0.4,-2.8); \node at (0.5,-3.0) {PLL (clocks)};
    \draw[FpgaOrange!70, line width=1pt] (0,-3.8) -- (0.4,-3.8); \node at (0.5,-3.8) {routing};
    \fill[FpgaOrange] (0.2,-4.5) circle (0.08); \node at (0.5,-4.5) {switch};
  \end{scope}
\end{tikzpicture}
\caption{A simplified layout of an FPGA chip}
\label{fig:fpga_grid}
\end{figure}

\section{The Main Idea: A Truth Table as Memory}

How can we build a universal logic element that can become AND, OR,
XOR, or anything else we like? Engineers came up with an elegant trick.

Any logic function of $n$ inputs is completely defined by its truth
table, which has $2^n$ rows. So it is enough to \emph{store} the right-hand
column of the table in a small memory and to use the inputs as the
\emph{address} of a cell. Such a device is called a \term{LUT} (Look-Up Table).

\begin{table}[H]
\centering
\caption{One LUT2 ``blank'' --- four different functions}
\label{tab:lut2}
\begin{tabular}{cc|c|cccc}
\toprule
$a$ & $b$ & address & AND & OR & NAND & XOR \\
\midrule
0 & 0 & 0 & 0 & 0 & 1 & 0 \\
0 & 1 & 1 & 0 & 1 & 1 & 1 \\
1 & 0 & 2 & 0 & 1 & 1 & 1 \\
1 & 1 & 3 & 1 & 1 & 0 & 0 \\
\midrule
\multicolumn{3}{r|}{\textbf{memory contents:}} & \code{1000} & \code{1110} & \code{0111} & \code{0110} \\
\bottomrule
\end{tabular}
\end{table}

To make a LUT ``become'' an AND gate, you write the numbers
$0,0,0,1$ into its memory; to make it an OR gate, you write $0,1,1,1$. No
rewiring is needed: only four configuration bits change.

\begin{figure}[H]
\centering
\begin{tikzpicture}[
  cell/.style={draw=FpgaDark, fill=FpgaLight, minimum width=9mm, minimum height=6.5mm, font=\ttfamily\small},
  mux/.style={draw=FpgaDark, fill=orange!15, thick, trapezium, trapezium angle=70,
              shape border rotate=270, minimum width=13mm, minimum height=6mm, font=\sffamily\tiny},
]
  % Memory cells
  \foreach \i/\v in {0/0,1/0,2/0,3/1} {
    \node[cell] (m\i) at (0,-\i*1.0) {\v};
    \node[lbl, left=1mm of m\i] {cell \i};
  }
  % Level 1 (controlled by b)
  \node[mux] (x0) at (2.6,-0.5) {MUX};
  \node[mux] (x1) at (2.6,-2.5) {MUX};
  \draw[wire] (m0.east) -- ++(0.5,0) |- ([yshift=1.3mm]x0.west);
  \draw[wire] (m1.east) -- ++(0.5,0) |- ([yshift=-1.3mm]x0.west);
  \draw[wire] (m2.east) -- ++(0.5,0) |- ([yshift=1.3mm]x1.west);
  \draw[wire] (m3.east) -- ++(0.5,0) |- ([yshift=-1.3mm]x1.west);
  % Level 2 (controlled by a)
  \node[mux] (y) at (5.0,-1.5) {MUX};
  \draw[wire] (x0.east) -- ++(0.4,0) |- ([yshift=1.3mm]y.west);
  \draw[wire] (x1.east) -- ++(0.4,0) |- ([yshift=-1.3mm]y.west);
  \draw[wire] (y.east) -- ++(1.0,0) node[right] {$y = f(a,b)$};
  % Control inputs
  \draw[wire, FpgaBlue] (2.6,-4.0) node[below] {$b$} -- (x1.south);
  \draw[wire, FpgaBlue] (2.6,-3.5) -| (1.9,-1.2) -| (x0.south);
  \draw[wire, FpgaBlue] (5.0,-4.0) node[below] {$a$} -- (y.south);
  % Frame
  \draw[dashed, FpgaGray, rounded corners] (-1.9,1.0) rectangle (5.7,-4.6);
  \node[lbl, anchor=south west] at (-1.9,1.0) {LUT2 configured as an AND gate};
  \node[lbl, anchor=south, align=center, text=FpgaGray] at (0,0.38) {configuration};
\end{tikzpicture}
\caption{The structure of a LUT2: four memory bits and a tree of multiplexers. The inputs
$a$ and $b$ select one of the cells}
\label{fig:lut2}
\end{figure}

The Gowin FPGA on the Tang Nano 20K uses
\term{LUT4}s --- lookup tables with 4 inputs. How many different functions can
a single LUT4 implement? Its memory holds 16 bits, and each can be 0 or 1:
\[
  N = 2^{2^4} = 2^{16} = 65\,536 \text{ functions.}
\]
These include AND, NOR, the majority gate, the half adder,
and in fact anything you can think of for four inputs.

\begin{figure}[H]
\centering
\begin{tikzpicture}
\begin{axis}[
  ybar, bar width=14mm, width=0.7\textwidth, height=6cm,
  ymode=log, log origin=infty, ymin=1, ymax=1e6,
  symbolic x coords={LUT1, LUT2, LUT3, LUT4},
  xtick=data, ylabel={Number of possible functions},
  nodes near coords, nodes near coords style={font=\sffamily\small},
  point meta=explicit symbolic,
  label style={font=\sffamily\small}, tick label style={font=\sffamily\small},
  grid=major, grid style={FpgaGray!20}, axis lines*=left,
]
  \addplot[fill=FpgaBlue!70, draw=FpgaBlue] coordinates {
    (LUT1,4) [4] (LUT2,16) [16] (LUT3,256) [256] (LUT4,65536) [65\,536]};
\end{axis}
\end{tikzpicture}
\caption{The number of functions a LUT with $n$ inputs can implement ($2^{2^n}$)}
\end{figure}

\begin{warnbox}[What if there are more than four inputs?]
A function of 5, 6, or more variables is split across several LUT4s
connected to each other. The synthesizer does this automatically, and
special MUX2 multiplexers next to the LUTs help combine two LUT4s into
one LUT5, two LUT5s into a LUT6, and so on. You will see them in the synthesis report
under the names \code{MUX2\_LUT5} and \code{MUX2\_LUT6}.
\end{warnbox}

\begin{diybox}[Find the table in a LUT]
\textbf{Task.} Find a LUT in the synthesized circuit and read the table stored in it.\\
\textbf{In Logra.} The ``Full adder'' example (New project $\to$ ``From an example''), ``Synthesize''; the ``Real board'' group, Gowin tab; click the carry LUT.\\
\textbf{Expected result.} In the ``Inspector'' there is a ``LUT with 3 inputs'' block and the INIT configuration \code{0xE8}, that is, the bits \code{11101000}: the ones are in rows 3, 5, 6, 7 --- where at least two of the inputs are ones (the majority function of the carry).\\
\textbf{How to check.} Expand the ``Truth table'' (8 rows) and compare it with Table~\ref{tab:lut2} using the rule ``address = inputs''. \stars{2} Read the INIT of the sum LUT and explain why its ones are in the rows with an odd number of ones.
\end{diybox}

\section{The Flip-Flop: The Circuit Gets Memory}

LUTs alone are not enough. From Chapter~\ref{ch:comb} you remember that there are
\emph{combinational} circuits (the output depends only on the current inputs) and
\emph{sequential} ones, which have memory. For memory, a
\term{D flip-flop} sits next to each LUT.

Recall (Chapter~\ref{ch:ff}): a D flip-flop remembers the value of input $D$ at
the moment of the \emph{rising edge} of the clock signal (the $0\to1$ transition) and holds
it on output $Q$ until the next edge.

\begin{figure}[H]
\centering
\begin{tikztimingtable}[timing/slope=0.1, timing/coldist=2pt, xscale=2.6, yscale=1.4,
  timing/font=\sffamily\small, timing/rowdist=3.2ex]
  CLK & 8{1L 1H} \\
  D   & 2L 4H 4L 2H 2L 2H \\
  Q   & 3L 4H 4L 2H 2L 1H \\
\extracode
  \begin{pgfonlayer}{background}
    \foreach \t in {1,3,...,15} \draw[FpgaOrange, dashed, thin] (\t,0.6) -- (\t,-5.0);
  \end{pgfonlayer}
  \node[lbl, text=FpgaOrange, anchor=south] at (8,1.6) {dashed lines --- rising edges of CLK};
\end{tikztimingtable}
\caption{Timing diagram of a D flip-flop: output $Q$ changes only on the
rising edge of CLK and takes the value that $D$ had at that moment}
\label{fig:dff}
\end{figure}

\begin{diybox}[The flip-flop catches the edge]
\textbf{Task.} Check on the model that a D flip-flop changes its output only on the rising edge.\\
\textbf{In Logra.} ``New project'' $\to$ ``From logic elements'': a D flip-flop, input $D$, a clock input, output $Q$; the ``Preview'' button.\\
\textbf{Expected result.} While the clock does not change, clicking $D$ has no effect on $Q$; $Q$ updates when the clock goes $0\to1$.\\
\textbf{How to check.} Compare the behavior with Fig.~\ref{fig:dff}: on the Waveforms tab set up a scenario (events on $D$ between edges) and make sure that $Q$ repeats $D$ only at the moments of the edge. \stars{2} Find the DFF element of this circuit on the Gowin tab.
\end{diybox}

\section{The CFU Logic Cell}

LUTs and flip-flops are combined into \term{logic cells}. In Gowin devices such a cell
is called a \term{CFU} (Configurable Function Unit). Inside a CFU there are four
``slices'' called CLS (Configurable Logic Slice); each has two LUT4s, and three of them
also have two flip-flops each. In total, one CFU has 8 LUT4s and 6 flip-flops.

Let us check: the GW2AR-18 FPGA contains $2592$ CFU cells, i.e.\
$2592\cdot 8 = 20\,736$ LUT4s and $2592\cdot 6 = 15\,552$ flip-flops. These numbers are
given in the chip's documentation.

\begin{figure}[H]
\centering
\begin{tikzpicture}[
  lut/.style={draw=FpgaBlue, fill=FpgaLight, thick, minimum width=16mm, minimum height=13mm, font=\sffamily\small, align=center},
  ff/.style={draw=FpgaGreen, fill=green!8, thick, minimum width=11mm, minimum height=13mm, font=\sffamily\small},
  mx/.style={draw=FpgaDark, fill=orange!15, thick, trapezium, trapezium angle=70, shape border rotate=270, minimum width=10mm, minimum height=5mm, font=\sffamily\tiny},
]
  \node[lut] (lut) at (0,0) {LUT4\\\scriptsize 16 bits};
  \foreach \i/\n in {1/A,2/B,3/C,4/D} {
    \draw[wire] ([yshift={(2.5-\i)*3.6mm}]lut.west) -- ++(-1,0) node[left, font=\sffamily\footnotesize, inner sep=1pt] {\n};
  }
  \node[ff] (ff) at (3.6,0) {D\quad Q};
  \draw[wire] (lut.east) -- node[above, lbl] {F} (ff.west);
  \draw[wire] ([xshift=-3mm]ff.south) -- ++(0,-0.8) node[below left=0mm, lbl] {CLK};
  \draw[wire] ([xshift=3mm]ff.south) -- ++(0,-0.8) node[below right=0mm, lbl] {CE, SR};
  \node[mx] (m) at (6.1,0) {MUX};
  \draw[wire] (ff.east) -- ([yshift=-1.2mm]m.west);
  \draw[wire] (1.6,0) |- (5.0,1.1) |- ([yshift=1.2mm]m.west);
  \fill (1.6,0) circle (1.3pt);
  \draw[wire] (m.east) -- ++(1,0) node[right, font=\sffamily\small] {Output};
  \draw[wire, FpgaGray] (m.south) -- ++(0,-0.7) node[below, lbl] {configuration bit};
  \node[lbl, align=center] at (3.1,1.45) {combinational output (bypassing the flip-flop)};
  \draw[dashed, FpgaGray, rounded corners] (-2.0,-1.9) rectangle (9.8,1.9);
  \node[lbl, anchor=north west] at (-2.0,1.9) {half of a CLS};
\end{tikzpicture}
\caption{The basic element of a logic cell: a LUT4, a D flip-flop, and a multiplexer
that selects whether the output comes from the flip-flop or directly from the LUT}
\label{fig:cls}
\end{figure}

\begin{diybox}[How many cells does a circuit occupy]
\textbf{Task.} Relate the resource counters to the structure of a CFU.\\
\textbf{In Logra.} The ``Running light'' project, ``Synthesize''; the ``Real board'' group, the Chip tab, ``Place and build''.\\
\textbf{Expected result.} The counters show LUT 15, FF 29, ALU 22; in the ``nextpnr result'' block the value ``Tiles with cells'' is 20.\\
\textbf{How to check.} One CFU contains 8 LUT4s and 6 flip-flops: estimate the minimum number of CFUs for 29 flip-flops ($\lceil 29/6\rceil$) and compare it with the number of occupied tiles; explain why more are occupied (flip-flops, LUTs, and ALUs are not always packed densely). \stars{2} Click an occupied tile on the map and read the names of its elements.
\end{diybox}

\section{Routing Resources}

Thousands of cells have to be connected to each other. For this purpose,
\term{routing channels} run between them --- bundles of wires of different lengths (short ones for
neighbors, long ones for distant blocks). At the crossings there are
\term{switch matrices}: transistor switches, each of which
is controlled by its own configuration bit.

\begin{figure}[H]
\centering
\begin{tikzpicture}
  \foreach \i in {0,1,2,3} {
    \draw[wire, FpgaGray] (\i*0.5,-0.5) -- (\i*0.5,2.5);
    \draw[wire, FpgaGray] (-0.5,\i*0.5+0.25) -- (2.0,\i*0.5+0.25);
  }
  \foreach \x in {0,1,2,3} \foreach \y in {0,1,2,3}
    \draw[FpgaGray] (\x*0.5,\y*0.5+0.25) circle (0.08);
  \foreach \x/\y in {0/1, 2/3, 3/0}
    \fill[FpgaOrange] (\x*0.5,\y*0.5+0.25) circle (0.1);
  \draw[very thick, FpgaOrange] (-0.5,0.75) -- (0,0.75) -- (0,2.5);
  \draw[very thick, FpgaOrange] (1.0,-0.5) -- (1.0,1.75) -- (2.0,1.75);
  \draw[very thick, FpgaOrange] (1.5,-0.5) -- (1.5,0.25) -- (2.0,0.25);
  \node[lbl, align=left, anchor=west] at (2.6,1.9) {\textcolor{FpgaOrange}{$\bullet$} switch closed (bit = 1)};
  \node[lbl, align=left, anchor=west] at (2.6,1.3) {$\circ$ switch open (bit = 0)};
  \node[lbl, align=left, anchor=west] at (2.6,0.5) {The orange lines are three\\ real ``wires'' of the circuit};
\end{tikzpicture}
\caption{A switch matrix: closed switches form the required connections}
\end{figure}

\begin{infobox}[A curious fact]
Most of the chip area of an FPGA is not occupied by logic at all, but by
routing and the configuration bits that control it. This is why FPGAs are
slower and more expensive than an ASIC of the same complexity: you pay for flexibility.
\end{infobox}

\begin{diybox}[Show the routing]
\textbf{Task.} See how a circuit uses the wires of the chip.\\
\textbf{In Logra.} The Chip tab after ``Place and build''; the ``Map layers'' panel on the right.\\
\textbf{Expected result.} Turn on the ``Routing'' layer (it is off by default): the darker a tile, the more wires in it are in use; in the ``nextpnr result'' block there is a ``Hops'' row (328).\\
\textbf{How to check.} Turn off the ``Logic'', ``Registers'', and ``Pins'' layers one by one and write down what disappears from the map. \stars{2} Compare the number of ``Hops'' with the number of placed elements (93) and draw a conclusion about which is more expensive: logic or wires (see ``A curious fact'').
\end{diybox}

\section{Specialized Blocks}

Besides general-purpose logic, an FPGA has ready-made ``bricks'' for common tasks.
You can build the same thing from LUTs, but it will be slower and take up a lot of space.

\begin{description}
  \item[IOB (input/output blocks)] connect the internal logic to the pins
    of the chip. In an IOB you can choose the voltage standard (LVCMOS 3.3~V,
    1.8~V, LVDS\ldots), the drive current, and a pull-up or pull-down.
  \item[BSRAM (block memory)] --- ready-made blocks of static memory of 18~Kbit each.
    The GW2AR-18 has 46 of them, 828~Kbit in total. They are used for buffers, FIFOs, and tables.
  \item[DSP (digital signal processing)] --- hardware multipliers
    of $18\times18$ bits with adders. There are 48 of them; they multiply in a single clock cycle.
  \item[PLL (phase-locked loop)] --- can produce almost any other frequency from the
    27~MHz input: 50, 100, 126, 371~MHz\ldots\
    The GW2AR-18 has two PLLs.
  \item[Global clock network] --- special wires that distribute
    the clock signal across the whole chip with minimal timing skew.
\end{description}

\begin{figure}[H]
\centering
\begin{tikzpicture}
\begin{axis}[
  xbar, bar width=4.5mm, width=0.78\textwidth, height=6.5cm,
  xmode=log, log origin=infty, xmin=1, xmax=2e5, xlabel={Count (logarithmic scale)},
  symbolic y coords={PLL, {18x18 multipliers}, {BSRAM blocks}, {Flip-flops (FF)}, {LUT4}},
  ytick=data, y tick label style={font=\sffamily\small},
  nodes near coords, nodes near coords style={font=\sffamily\small, anchor=west},
  point meta=explicit symbolic,
  label style={font=\sffamily\small}, x tick label style={font=\small, /pgf/number format/1000 sep={\,}},
  axis lines*=left, grid=major, grid style={FpgaGray!20}, enlarge y limits=0.15,
]
  \addplot[fill=FpgaBlue!70, draw=FpgaBlue] coordinates {
    (20736,{LUT4}) [20\,736] (15552,{Flip-flops (FF)}) [15\,552]
    (46,{BSRAM blocks}) [46] (48,{18x18 multipliers}) [48] (2,PLL) [2]};
\end{axis}
\end{tikzpicture}
\caption{Resources of the GW2AR-18 FPGA on the Tang Nano 20K board}
\end{figure}

\begin{diybox}[Multiplier and memory]
\textbf{Task.} Find out which special blocks are used by a multiplier and by memory.\\
\textbf{In Logra.} ``Designer'': place a multiplier (then, in a separate project, a RAM), ``Synthesize''; the Chip tab, the BRAM and DSP counters.\\
\textbf{Expected result.} For the ``Running light'' the BRAM and DSP counters are 0; in the new projects your own values are shown (a multiplier can be built from a DSP block or from LUTs and ALUs, a RAM from BSRAM or from flip-flops).\\
\textbf{How to check.} Turn on the ``Memory and DSP'' layer and see whether colored tiles appear; compare with the GW2AR-18 counters from this chapter (46 BSRAM, 48 DSP).
\end{diybox}

\section{Where the Configuration Is Stored}

The configuration memory of a Gowin FPGA is built with SRAM technology, like the
RAM of a computer: it is erased when the power is turned off. That is why
an FPGA has two loading modes:
\begin{itemize}
  \item \textbf{into SRAM} --- fast and convenient for debugging, but only until the power is switched off;
  \item \textbf{into Flash} --- the configuration file is written to flash memory,
        and every time it is powered on, the FPGA loads its circuit from there by itself
        within tens of milliseconds.
\end{itemize}

\begin{ideabox}
Inside an FPGA: \textbf{LUTs} (any logic function), \textbf{flip-flops}
(1 bit of memory), \textbf{routing} (wires and switches), \textbf{IOBs}
(pins), as well as ready-made blocks of \textbf{memory}, \textbf{multipliers}, and
\textbf{PLLs}. The configuration is millions of bits that define
the contents of every LUT and the state of every switch.
\end{ideabox}

\begin{diybox}[The configuration file]
\textbf{Task.} Find the configuration file and learn how it gets into the chip.\\
\textbf{In Logra.} The Chip tab: the ``Bitstream .fs'' link after the build; then the Connect tab: three steps and the ``Load into SRAM'' button.\\
\textbf{Expected result.} The \code{.fs} file is downloaded; the Connect tab shows that loading goes into SRAM, that is, only until the power is switched off.\\
\textbf{How to check.} Write down the file size and repeat for the ``27 MHz blinker'' project: the sizes are usually close, because the configuration describes the whole chip, not just the occupied part. \stars{2} Explain why, after the power is cut, the circuit has to be loaded again unless it was written to Flash.
\end{diybox}
